% Source: https://artofproblemsolving.com/community/c3h3146489_2023_bmmt_team_round__berkley_mini_math_tournament_fall
% Stable identity: aops-topic-3146489-q13
If adjacent terms had the same units digit, an even digit would give common factor $2$, a digit $5$ would give common factor $5$, and an odd digit other than $5$ would force both numbers to be divisible by $3$. Thus only the units digit can change, and all terms share the same first two digits. Write them as $b+d$, where $10\mid b$ and $0\le d\le9$, with increasing $d$. If $b\equiv0\pmod3$, the allowed digits are $0,2,3,4,5,6,8,9$. All eight cannot be used, since the first two would both be even. If $b\equiv1\pmod3$, only $0,2,4,5,6,8$ are allowed, at most six. If $b\equiv2\pmod3$, the allowed digits are $0,1,2,4,5,6,7,8$; using all eight would place the even-ending digits $2,4$ adjacent. Hence the length is at most seven. The sequence $110,111,112,115,116,117,118$ achieves seven: each term satisfies the divisibility condition, each consecutive pair is coprime, and only the units digit changes. Thus the maximum is $\boxed7$.
